%%%%%%%%%%%%%%%%%%%%%%%%%%%%%%%%%%%%%%%%%%%%%%%%%%%%%%%
% Arquivo para entrada de dados para a parte pré textual
%%%%%%%%%%%%%%%%%%%%%%%%%%%%%%%%%%%%%%%%%%%%%%%%%%%%%%%
% 
% Basta digitar as informações indicidas, no formato 
% apresentado.
%
%%%%%%%
% Os dados solicitados são, na ordem:
%
% tipo do trabalho
% componentes do trabalho 
% título do trabalho
% nome do autor
% local 
% data (ano com 4 dígitos)
% orientador(a)
% coorientador(a)(as)(es)
% arquivo com dados bibliográficos
% instituição
% setor
% programa de pós gradução
% curso
% preambulo
% data defesa
% CDU
% errata
% assinaturas - termo de aprovação
% resumos & palavras chave
% agradecimentos
% dedicatoria
% epígrafe


% ---
% Informações de dados para CAPA e FOLHA DE ROSTO
% ---
\titulo{O Efeito das Eleições nas Finanças Municipais: Evidências Empíricas para os municípios brasileiros}
\autor{Lorenzo Rodrigues Geronazzo}
\local{Curitiba}
\data{2026}
\orientador{Victor Rodrigues de Oliveira}

\tipotrabalho{Monografia}
% O preambulo deve conter o tipo do trabalho, o objetivo, 
% o nome da instituição e a área de concentração 
% Sem ponto final: o modelo acrescenta "." na folha de rosto e ", pela seguinte
% banca examinadora:" no termo de aprovação.
\preambulo{Monografia apresentada ao curso de Ciências Econômicas, Setor de Ciências Sociais Aplicadas, da Universidade Federal do Paraná, como requisito parcial à obtenção do título de Bacharel em Ciências Econômicas}

% Marcar Sim para as partes que irão compor o documento pdf
%----------------------------------------------------------------------------- 
 \providecommand{\terCapa}{Sim}
 \providecommand{\terFolhaRosto}{Sim}
 \providecommand{\terTermoAprovacao}{Sim}
 \providecommand{\terDedicatoria}{Sim}
 \providecommand{\terFichaCatalografica}{Nao}
 \providecommand{\terEpigrafe}{Sim}
 \providecommand{\terAgradecimentos}{Sim}
 \providecommand{\terErrata}{Nao}
 \providecommand{\terListaFiguras}{Sim}
 \providecommand{\terListaTabelas}{Sim}
 \providecommand{\terSiglasAbrev}{Sim}
 \providecommand{\terResumos}{Sim}
 \providecommand{\terSumario}{Sim}
 \providecommand{\terAnexo}{Nao}
 \providecommand{\terApendice}{Sim}
 \providecommand{\terIndiceR}{Nao}
%----------------------------------------------------------------------------- 

\titulo{O Efeito das Eleições nas Finanças Municipais: Evidências Empíricas para os municípios brasileiros}
\autor{Lorenzo Rodrigues Geronazzo}
\local{Curitiba}
\data{2026} %Apenas ano 4 dígitos

% Orientador ou Orientadora
\orientador{Victor Rodrigues de Oliveira}
% \orientadora{Nome da sua Orientadora Aqui}
% Se houver os dois prevalece o feminino.

% Em termos de coorientação, podem haver até quatro neste modelo
% Sendo 2 mulhere e 2 homens.
% Coorientador ou Coorientadora
%\coorientador{}%Prof Morgan Freeman, DSc}
%\coorientadora{Prof\textordfeminine~Audrey Hepburn, DEng}

% Segundo Coorientador ou Segunda Coorientadora
\scoorientador{}
%Prof Jack Nicholson, DEng}
\scoorientadora{}
%Prof\textordfeminine~Ingrid Bergman, DEng}
% ----------------------------------------------------------
\addbibresource{referencias.bib}

% ----------------------------------------------------------
\instituicao{%
Universidade Federal do Paraná}

\def \ImprimirSetor{Setor de Ciências Sociais Aplicadas}

% Monografia de graduação: sem programa de pós-graduação
\def \ImprimirProgramaPos{}

\def \ImprimirCurso{Curso de Ciências Econômicas}

% (o \preambulo fica no início deste arquivo, junto do \tipotrabalho)

%-----------------------------------------------------------------------------

\newcommand{\imprimirCurso}{Ciências Econômicas}

\newcommand{\imprimirDataDefesa}{[a definir]}

\newcommand{\imprimircdu}{[a definir]}

% ----------------------------------------------------------
\newcommand{\imprimirerrata}{
Elemento opcional da \cites[4.2.1.2]{NBR14724:2011}. Exemplo:

\vspace{\onelineskip}

FERRIGNO, C. R. A. \textbf{Tratamento de neoplasias ósseas apendiculares com
reimplantação de enxerto ósseo autólogo autoclavado associado ao plasma
rico em plaquetas}: estudo crítico na cirurgia de preservação de membro em
cães. 2011. 128 f. Tese (Livre-Docência) - Faculdade de Medicina Veterinária e
Zootecnia, Universidade de São Paulo, São Paulo, 2011.

\begin{table}[htb]
\center
\footnotesize
\begin{tabular}{|p{1.4cm}|p{1cm}|p{3cm}|p{3cm}|}
  \hline
   \textbf{Folha} & \textbf{Linha}  & \textbf{Onde se lê}  & \textbf{Leia-se}  \\
    \hline
    1 & 10 & auto-conclavo & autoconclavo\\
   \hline
\end{tabular}
\end{table}}

% Comandos de dados - Data da apresentação
\providecommand{\imprimirdataapresentacaoRotulo}{}
\providecommand{\imprimirdataapresentacao}{}
\newcommand{\dataapresentacao}[2][\dataapresentacaoname]{\renewcommand{\dataapresentacao}{#2}}

% Comandos de dados - Nome do Curso
\providecommand{\imprimirnomedocursoRotulo}{}
\providecommand{\imprimirnomedocurso}{}
\newcommand{\nomedocurso}[2][\nomedocursoname]
  {\renewcommand{\imprimirnomedocursoRotulo}{#1}
\renewcommand{\imprimirnomedocurso}{#2}}


% ----------------------------------------------------------
\newcommand{\AssinaAprovacao}{

% Banca examinadora (o orientador é impresso antes, pelo modelo).
% \mbox: sem ele, o "[" logo depois do \\ interno do \assinatura vira argumento opcional.
\assinatura{\mbox{[a definir]}}
   \assinatura{\mbox{[a definir]}}
      
   \begin{center}
    \vspace*{0.5cm}
    %{\large\imprimirlocal}
    %\par
    %{\large\imprimirdata}
    \imprimirlocal, \imprimirDataDefesa.
    \vspace*{1cm}
  \end{center}
  }
  
% ----------------------------------------------------------
%\newcommand{\Errata}{%\color{blue}
%Elemento opcional da \textcite[4.2.1.2]{NBR14724:2011}. Exemplo:
%}

% ----------------------------------------------------------
\newcommand{\EpigrafeTexto}{%\color{blue}
\textit{``Não vos amoldeis às estruturas deste mundo, \\
		mas transformai-vos pela renovação da mente, \\
		a fim de distinguir qual é a vontade de Deus: \\
		o que é bom, o que Lhe é agradável, o que é perfeito.\\
		(Bíblia Sagrada, Romanos 12, 2)}
}

% ----------------------------------------------------------
\newcommand{\ResumoTexto}{%\color{blue}
Este trabalho investiga como as despesas dos municípios brasileiros respondem ao calendário eleitoral e ao alinhamento político do prefeito com o governador e com o presidente, entre 2013 e 2025. Com dados do Sistema de Informações Contábeis e Fiscais do Setor Público Brasileiro (Siconfi), do Tribunal Superior Eleitoral e do Instituto Brasileiro de Geografia e Estatística (IBGE) para cerca de 5,6 mil municípios, as despesas são agrupadas por função de governo e analisadas em dois modelos de efeitos fixos: um para o ciclo eleitoral, com erros-padrão de Driscoll e Kraay, e outro para o alinhamento político, com efeitos fixos de município e de ano. A despesa entra em logaritmo, e a robustez dos resultados é avaliada em dezoito testes. Os resultados indicam que o ciclo eleitoral aparece sobretudo na composição do gasto: no ano eleitoral, as despesas com Habitação e Urbanismo, Transporte e Saúde e Saneamento são, em média, cerca de 27\%, 19,5\% e 4,2\% maiores, com aumentos nas duas eleições analisadas, enquanto Administração e Educação e Cultura perdem participação na despesa total. A despesa total é cerca de 5\% maior no ano eleitoral, efeito mais forte em 2024 e sensível ao período da amostra. O alinhamento com o governador está associado a uma despesa total cerca de 0,4\% maior, e o alinhamento com o presidente não tem efeito sobre a despesa total na amostra completa. Os resultados são compatíveis com uma recomposição do gasto em favor das despesas mais visíveis ao eleitor e confirmam, para um período recente, parte das evidências encontradas para os municípios brasileiros entre 1990 e 2005.
}

\newcommand{\PalavraschaveTexto}{%\color{blue}
ciclos políticos orçamentários; finanças municipais; dados em painel; alinhamento político.}

% ----------------------------------------------------------
\newcommand{\AbstractTexto}{%\color{blue}
This study investigates how Brazilian municipal expenditures respond to the electoral calendar and to the political alignment of mayors with state governors and the president between 2013 and 2025. Using data from the Brazilian Public Sector Accounting and Tax Information System (Siconfi), the Superior Electoral Court and the Brazilian Institute of Geography and Statistics (IBGE) for about 5,600 municipalities, expenditures are grouped by government function and analyzed in two fixed-effects models: one for the electoral cycle, with Driscoll-Kraay standard errors, and another for political alignment, with municipality and year fixed effects. Expenditure enters in logarithms, and the robustness of the results is assessed in eighteen tests. The results indicate that the electoral cycle appears mainly in the composition of spending: in election years, spending on Housing and Urban Planning, Transportation, and Health and Sanitation is, on average, about 27\%, 19.5\% and 4.2\% higher, respectively, with increases in both elections analyzed, while Administration and Education and Culture lose share of total spending. Total spending is about 5\% higher in election years, an effect that is stronger in 2024 and sensitive to the sample period. Alignment with the governor is associated with total spending about 0.4\% higher, while alignment with the president has no effect on total spending in the full sample. The results are consistent with a reallocation of spending toward items more visible to voters and confirm, for a recent period, part of the evidence found for Brazilian municipalities between 1990 and 2005.
}
% ---
\newcommand{\KeywordsTexto}{%\color{blue}
political budget cycles; municipal finance; panel data; political alignment.
}

% ----------------------------------------------------------
\newcommand{\Resume}
{%\color{blue}
%Il s'agit d'un résumé en français.
} 
% ---
\newcommand{\Motscles}
{%\color{blue}
 %latex. abntex. publication de textes.
}

% ----------------------------------------------------------
\newcommand{\Resumen}
{%\color{blue}
%Este es el resumen en español.
}
% ---
\newcommand{\Palabrasclave}
{%\color{blue}
%latex. abntex. publicación de textos.
}

% ----------------------------------------------------------
\newcommand{\AgradecimentosTexto}{%
Agradeço ao meu orientador, Prof. Victor Rodrigues de Oliveira, pela orientação, pela disponibilidade e pelas revisões cuidadosas ao longo deste trabalho.

À Universidade Federal do Paraná e aos professores do curso de Ciências Econômicas, pela formação que tornou este trabalho possível.

Aos meus pais e à minha família, pelo apoio e incentivo durante toda a graduação.
}

% ----------------------------------------------------------
\newcommand{\DedicatoriaTexto}{%
Dedico este trabalho aos meus pais, Ricardo Luiz Geronazzo e Denise Maria Rodrigues Geronazzo, que sempre me incentivaram a estudar e a explorar meu potencial; à minha namorada, Pietra Pizzatto Minatti, pela paciência e pelo apoio ao longo desta etapa; ao meu tio, Renê Renato Rodrigues, por acreditar em mim; e ao meu avô, Lourenço José Geronaso, pelo exemplo.
}

